\section*{How the Course Works}
\addcontentsline{toc}{section}{How the Course Works}

\begin{corebox}[Two uses of the same chapter]
Each chapter supports the live lecture and later independent study. During the
lecture it provides the mathematical route and the core results to discuss.
After the lecture it remains a complete narrative containing definitions,
examples, arguments, and interpretations that the student can revisit without
the lecturer being present.
\end{corebox}

This book is part of a weekly learning process rather than material to be read
only before an examination. Each of the twelve thematic chapters is designed to
support one week of study. In a fifteen-week semester, the remaining time can be
used to introduce the method of work, connect topics, revise important ideas,
and present more substantial solutions.

\subsection*{The Weekly Learning Cycle}

\begin{corebox}[The complete weekly cycle]
\begin{enumerate}
    \item \textbf{Attend the lecture.}
    The lecture introduces the topic, motivates the central questions, and
    demonstrates how the relevant objects and methods should be read.

    \item \textbf{Return to the chapter.}
    Read the corresponding chapter carefully after the lecture. The PDF contains
    the full argument, not merely the material shown or discussed in detail live.

    \item \textbf{Ask questions and clarify the material.}
    Mark every unclear definition, transition, calculation, or conclusion and
    resolve it using the text, discussion, examples, and appropriate tools.

    \item \textbf{Solve the complete assigned problem set.}
    Move from following explanations to constructing arguments independently.

    \item \textbf{Prepare professional mathematical notes.}
    Use precise notation, explanatory sentences, justified transitions, and a
    clear conclusion.

    \item \textbf{Present and discuss your work.}
    Be prepared to explain a solution, answer questions about its steps, and
    describe what was learned while preparing it.
\end{enumerate}
\end{corebox}

\begin{interpretationbox}[Why the stages form one process]
The lecture prepares the reading, the reading prepares the problem solving, the
written solution prepares the presentation, and feedback leads back to revision.
No stage is intended to replace the others.
\end{interpretationbox}

\subsection*{Exercise Classes: Presentation, Feedback, and Revision}

\begin{corebox}[What a presentation must demonstrate]
A student should be able to read the notation, explain the role of each important
formula, justify the method, and respond to questions without relying on AI to
speak on the student's behalf.
\end{corebox}

Presenting mathematics is part of learning mathematics. It reveals whether an
argument has genuinely been understood and whether it can be communicated
clearly. Feedback given during exercise classes is intended for the whole group,
not only for the student whose work is being presented. By comparing different
solutions, students learn which explanations are complete, which steps need
justification, and which weaknesses require revision.

\begin{remarkbox}
Feedback concerns the mathematical work and its results, not the personal worth
of the student. It should identify what already works and what must be changed to
produce a stronger argument.
\end{remarkbox}

\begin{interpretationbox}[Revision is part of mathematical work]
The first submitted version need not be the final version. Receiving comments,
reconsidering choices, correcting errors, and adding missing explanations are
normal parts of serious academic and professional work, not evidence of failure.
\end{interpretationbox}

\subsection*{Progress and Assessment}

\begin{corebox}[What progress looks like]
By the end of a weekly cycle, the student should be able to identify the main
objects, state the essential definitions, solve representative problems, justify
the principal steps, verify or interpret the result, and communicate the argument
clearly.
\end{corebox}

Evidence of serious engagement includes:
\begin{itemize}
    \item regular preparation and completion of assigned problem sets;
    \item documented attempts rather than bare answers;
    \item active participation in presentations and discussion;
    \item separate, self-contained notes for individual problems;
    \item careful revision after feedback;
    \item increasing clarity, independence, and mathematical precision.
\end{itemize}

\begin{remarkbox}
A weak first solution or presentation does not determine the final exercise
component grade. Sustained work, revision, and demonstrated development can lead
to a very good final result. This does not lower the standard; it recognises the
process required to reach it.
\end{remarkbox}

The examination is assessed separately and tests mathematical knowledge and the
ability to use it independently. Regular work does not replace the examination,
but it develops the understanding and skills needed to complete it successfully.

\begin{summarybox}[How to use the course]
Attend, reread, question, solve, write, present, and revise. The lecture provides
orientation; the chapter preserves the complete argument; independent work turns
that argument into understanding.
\end{summarybox}
